% SEGMENT OLP-0070-S01
% Part: first-order-logic
% Chapter: sequent-calculus
% Section: rules-and-proofs

\documentclass[../../../include/open-logic-section]{subfiles}

\begin{document}

\iftag{FOL}
      {\olfileid{fol}{seq}{rul}}
      {\olfileid{pl}{seq}{rul}}

\olsection{விதிகள் மற்றும் \usetoken{P}{derivation}}

பின்வருவனவற்றில், $\Gamma, \Delta, \Pi, \Lambda$ ஆகியவை
!!{sentence}s[gen] முடிவுறு தொடர்களைக் குறிப்பனவாக இருக்கட்டும்.

\begin{defn}[தொடரணி]
\emph{தொடரணி} என்பது
\[
\Gamma \Sequent \Delta
\]
என்ற வடிவிலான கோவை; இங்கு $\Gamma$, $\Delta$ ஆகியவை மொழி $\Lang L$ இன்
!!{sentence}s[acc] கொண்ட முடிவுறு (வெறுமையாகவும் இருக்கக்கூடிய) தொடர்கள்.
$\Gamma$ \emph{முன்தொடர்} என்றும், $\Delta$ \emph{பின்தொடர்} என்றும்
அழைக்கப்படுகின்றன.
\end{defn}

% SEGMENT OLP-0070-S02
\begin{explain}
ஒரு தொடரணியின் உள்ளுணர்வுக் கருத்து இதுவே: முன்தொடரிலுள்ள !!{sentence}s
அனைத்தும் மெய்யெனில், பின்தொடரிலுள்ள !!{sentence}s[loc] குறைந்தது ஒன்றாவது
மெய். அதாவது, $\Gamma = \tuple{!A_1, \dots, !A_m}$ என்றும்
$\Delta = \tuple{!B_1, \dots, !B_n}$ என்றும் இருந்தால், $\Gamma \Sequent
\Delta$ மெய்யாவது, அப்போதும் அப்போது மட்டுமே,
\[
(!A_1 \land \cdots \land !A_m) \lif (!B_1 \lor \cdots \lor
!B_n)
\]
மெய். இரண்டு சிறப்பு நேர்வுகள் உள்ளன: $\Gamma$ வெறுமையாக இருப்பதும்
$\Delta$ வெறுமையாக இருப்பதும். $\Gamma$ வெறுமையாக, அதாவது $m = 0$ ஆக
இருந்தால், $\quad \Sequent \Delta$ மெய்யாவது, அப்போதும் அப்போது மட்டுமே,
$!B_1 \lor \dots \lor !B_n$ மெய். $\Delta$ வெறுமையாக, அதாவது $n = 0$
ஆக இருந்தால், $\Gamma \Sequent \quad$ மெய்யாவது, அப்போதும் அப்போது மட்டுமே,
$\lnot(!A_1 \land \dots \land !A_m)$ மெய். தொடர்புடைய !!{sentence}
செல்லுபடியாக இருந்தால், அப்போதும் அப்போது மட்டுமே, அந்தத் தொடரணி
செல்லுபடியாகும் என்கிறோம்.
\end{explain}

% SEGMENT OLP-0070-S03
$\Gamma$ என்பது !!{sentence}s[gen] ஒரு தொடர் எனில், $\Gamma$ தொடரின் வலது முனையில்
$!A$ ஐச் சேர்ப்பதால் கிடைப்பதை $\Gamma, !A$ என்றும் ($\Gamma$ தொடரின் இடது முனையில் $!A$
ஐச் சேர்ப்பதால் கிடைப்பதை $!A, \Gamma$ என்றும்) எழுதுகிறோம். $\Delta$வும்
!!{sentence}s[gen] ஒரு தொடர் எனில், $\Gamma, \Delta$ என்பது அவ்விரு
தொடர்களின் இணைத்தொடர்.

% SEGMENT OLP-0070-S04
\begin{defn}[தொடக்கத் தொடரணி]
\emph{தொடக்கத் தொடரணி} என்பது எந்த !!{sentence} $!A$ க்கும்
\iftag{prvFalse,prvTrue}{பின்வரும் வடிவங்களில் ஒன்றிலுள்ள தொடரணி:
  \begin{enumerate}
    \item $!A \Sequent !A$
    \tagitem{prvTrue}{$\quad \Sequent \ltrue$}{}
    \tagitem{prvFalse}{$\lfalse \Sequent \quad$}{}
  \end{enumerate}}
{$!A \Sequent !A$ என்ற வடிவிலுள்ள தொடரணி}.
\end{defn}

% SEGMENT OLP-0070-S05
தொடரணிக் கணியத்தின் !!^{derivation}s தொடரணிகளின் சில குறிப்பிட்ட மரங்கள்.
அவற்றின் உச்சியிலுள்ள தொடரணிகள் தொடக்கத் தொடரணிகளாக இருக்கும்; ஒரு தொடரணி
வேறு ஒன்று அல்லது இரண்டு தொடரணிகளுக்குக் கீழே இருந்தால், அது ஓர் அனுமான
விதியால் அவற்றிலிருந்து சரியாகப் பின்வர வேண்டும். $\Log{LK}$ க்கான விதிகள்
இரு முதன்மை வகைகளாகப் பிரிக்கப்படுகின்றன: \emph{தருக்க} விதிகள் மற்றும்
\emph{கட்டமைப்பு} விதிகள். கீழ்த் தொடரணியில் $!A$ மற்றும்/அல்லது $!B$ ஐக்
கொண்ட !!{sentence} ஒன்றுக்குரிய !!{main
  operator} பெயரால் தருக்க விதிகள்
அழைக்கப்படுகின்றன. ஒவ்வொரு விதிக்கும் இரண்டு வடிவங்கள் உண்டு: !!{operator}
கொண்ட !!{sentence} இடப்புறத்தில் உள்ள தொடரணியை அனுமானிக்கும் ஒரு வடிவமும்,
அந்த !!{sentence} வலப்புறத்தில் உள்ள தொடரணியை அனுமானிக்கும் ஒரு வடிவமும்.

\end{document}
